\documentclass[11pt,a4paper]{article}

% ---------------------------------------------------------------------
%  Paquetes
% ---------------------------------------------------------------------
% Se compila con LuaLaTeX para usar las mismas fuentes que la web (Fredoka y JetBrains Mono)
\usepackage{fontspec}
\usepackage[spanish,es-noshorthands,es-tabla]{babel}
\usepackage{microtype}
\usepackage[margin=2.5cm]{geometry}
\usepackage[dvipsnames,table]{xcolor}
\usepackage{amssymb}
\usepackage{float}
\usepackage{eso-pic}
\usepackage{pagecolor}
\usepackage{graphicx}
\usepackage{booktabs}
\usepackage{tabularx}
\usepackage{longtable}
\usepackage{array}
\usepackage{enumitem}
\usepackage{fancyhdr}
\usepackage{titlesec}
\usepackage{listings}
\usepackage[most]{tcolorbox}
\usepackage{tikz}
\usetikzlibrary{positioning,arrows.meta,shapes.geometric}
\usepackage{csquotes}
\usepackage[backend=biber,style=ieee]{biblatex}
\addbibresource{bibliografia.bib}
\usepackage[colorlinks]{hyperref}
\usepackage{xurl}

% =====================================================================
%  DATOS PERSONALES Y DE LA ACTIVIDAD  (rellenar antes de compilar)
% =====================================================================
\newcommand{\asignatura}{Programación y Bases de Datos}
\newcommand{\actividad}{Diseño de un recurso educativo gamificado: \emph{Misión SQL}}
\newcommand{\bloque}{Bloque temático 3: Bases de datos}
\newcommand{\materiaFP}{Módulo Profesional 0484 \emph{Bases de datos} (1.\textsuperscript{er} curso de los CFGS de Desarrollo de Aplicaciones Multiplataforma y Desarrollo de Aplicaciones Web)}

\newcommand{\alumno}{José Antonio Hernández López}
\newcommand{\dni}{48835386S}
\newcommand{\email}{jhernande2931@alumno.uned.es}
\newcommand{\curso}{Curso 2026--2027}

% --- URL públicas de los recursos (sustituir cuando estén publicados) ---
\newcommand{\urlGenially}{https://view.genially.com/XXXXXXXX}
% Los tres juegos son páginas web propias (HTML + CSS + JavaScript)
% publicadas con GitHub Pages en un único repositorio.
\newcommand{\urlRepo}{https://github.com/Antolin1/programacion-bbdd-uned}
\newcommand{\urlJuegos}{https://antolin1.github.io/programacion-bbdd-uned/}
\newcommand{\urlJuegoUno}{https://antolin1.github.io/programacion-bbdd-uned/juego1/}
\newcommand{\urlJuegoDos}{https://antolin1.github.io/programacion-bbdd-uned/juego2/}
\newcommand{\urlJuegoTres}{https://antolin1.github.io/programacion-bbdd-uned/juego3/}

% --- Número de diapositivas de la presentación Genially ---
\newcommand{\numDiapositivas}{[XX]}


% Nombres de \autoref en minúscula (se usan a mitad de frase: «la \autoref{tab:ra}» -> «la tabla 2»)
\addto\extrasspanish{%
  \def\tableautorefname{tabla}%
  \def\figureautorefname{figura}%
  \def\sectionautorefname{apartado}%
  \def\subsectionautorefname{apartado}%
  \def\subsubsectionautorefname{apartado}%
  \def\appendixautorefname{apéndice}%
  \def\equationautorefname{ecuación}}

\hypersetup{
  linkcolor=tierra, citecolor=hoja, urlcolor=acequia,
  pdftitle={Misión SQL -- Memoria de la actividad},
  pdfauthor={\alumno},
  pdfsubject={\asignatura}
}

% ---------------------------------------------------------------------
%  Estilo
% ---------------------------------------------------------------------
% --- Fuentes de la web (Fredoka no tiene cursiva: se inclina la redonda) ---
\setmainfont{Fredoka}[
  Path=fuentes/, Extension=.ttf,
  UprightFont=*-Regular, BoldFont=*-SemiBold,
  ItalicFont=*-Regular, ItalicFeatures={FakeSlant=0.18},
  BoldItalicFont=*-SemiBold, BoldItalicFeatures={FakeSlant=0.18}]
\newfontfamily\titular{Fredoka}[Path=fuentes/, Extension=.ttf, UprightFont=*-Bold]
\setmonofont{JetBrains Mono}[Scale=0.86, Ligatures=NoCommon, Contextuals=NoAlternate]

% --- Paleta de la huerta murciana (la misma que web/assets/common.css) ---
\definecolor{tierra}{HTML}{3A2A17}    % texto
\definecolor{tenue}{HTML}{7A6850}     % texto secundario
\definecolor{crema}{HTML}{F6EEDB}     % fondo
\definecolor{papel}{HTML}{FFFAF0}     % paneles
\definecolor{borde}{HTML}{DCCDB0}
\definecolor{limon}{HTML}{F4C430}
\definecolor{limontinta}{HTML}{946800}
\definecolor{naranja}{HTML}{D9691C}
\definecolor{verde}{HTML}{3F8A32}
\definecolor{hoja}{HTML}{2F5D24}
\definecolor{pimiento}{HTML}{C4432C}
\definecolor{acequia}{HTML}{2A7AA3}
\definecolor{berenjena}{HTML}{7B3F73}
% Nombres que usa el resto del documento
\colorlet{principal}{hoja}
\colorlet{acento}{naranja}
\definecolor{fondocodigo}{HTML}{FFFDF7}
\newcommand{\trab}{\textcolor{verde}{\ensuremath{\checkmark}}}

\pagecolor{papel}
\color{tierra}

% Franja a rayas de limón y hoja en lo alto de cada página, como la cabecera de la web
\newcommand{\toldo}[1]{%
  \begin{tikzpicture}[remember picture,overlay]
    \foreach \i in {0,...,30} {
      \fill[limon] ([xshift=\i*0.88cm, yshift=-#1]current page.north west) rectangle ++(0.44cm,-0.16cm);
      \fill[verde] ([xshift=\i*0.88cm+0.44cm, yshift=-#1]current page.north west) rectangle ++(0.44cm,-0.16cm);
    }
  \end{tikzpicture}}

\titleformat{\section}{\titular\LARGE\color{tierra}}{\textcolor{naranja}{\thesection}}{0.5em}{}
\titleformat{\subsection}{\titular\Large\color{hoja}}{\textcolor{naranja}{\thesubsection}}{0.5em}{}
\titleformat{\subsubsection}{\titular\large\color{tierra}}{\textcolor{naranja}{\thesubsubsection}}{0.5em}{}

\pagestyle{fancy}
\fancyhf{}
\fancyhead[L]{\small\color{naranja}\textbf{MISIÓN SQL}}
\fancyhead[R]{\small\color{tenue}\asignatura}
\fancyfoot[C]{\small\color{tenue}\thepage}
\renewcommand{\headrulewidth}{0pt}
\AddToShipoutPictureBG{\ifnum\value{page}>1 \toldo{0cm}\fi}

\setlist{itemsep=2pt,topsep=4pt}
\renewcommand{\arraystretch}{1.25}

\lstdefinestyle{sql}{
  language=SQL,
  basicstyle=\ttfamily\small,
  keywordstyle=\color{acequia}\bfseries,
  commentstyle=\color{tenue}\itshape,
  stringstyle=\color{pimiento},
  numberstyle=\color{tenue},
  backgroundcolor=\color{fondocodigo},
  frame=single, rulecolor=\color{borde},
  framesep=4pt, xleftmargin=6pt, xrightmargin=6pt,
  breaklines=true, columns=fullflexible, keepspaces=true,
  showstringspaces=false, upquote=true,
  morekeywords={REFERENCES,START,TRANSACTION,COMMIT,ROLLBACK,SAVEPOINT,
                TRUNCATE,GRANT,REVOKE,LIMIT,IF,EXISTS,CHECK},
  literate={á}{{\'a}}1 {é}{{\'e}}1 {í}{{\'i}}1 {ó}{{\'o}}1 {ú}{{\'u}}1
           {Á}{{\'A}}1 {É}{{\'E}}1 {Í}{{\'I}}1 {Ó}{{\'O}}1 {Ú}{{\'U}}1
           {ñ}{{\~n}}1 {Ñ}{{\~N}}1 {ü}{{\"u}}1 {º}{{\textordmasculine}}1
}
\lstset{style=sql}

% Cajas con el aspecto de los paneles de la web: fondo claro, borde suave y esquinas redondeadas
\newtcolorbox{reto}[1]{enhanced, breakable, colback=fondocodigo, colframe=acequia,
  coltitle=white, fonttitle=\titular, title={#1}, arc=3mm, boxrule=0.8pt,
  drop fuzzy shadow=borde}
\newtcolorbox{solucion}{enhanced, breakable, colback=verde!6!papel, colframe=verde,
  coltitle=white, fonttitle=\titular, title={Solución}, arc=3mm, boxrule=0.6pt}
\newtcolorbox{nota}{enhanced, colback=limon!16!papel, colframe=limon!50!naranja,
  arc=3mm, boxrule=0pt, borderline west={3pt}{0pt}{naranja}}

% =====================================================================
\begin{document}
% =====================================================================

% ---------------------------------------------------------------------
%  PORTADA
% ---------------------------------------------------------------------
\begin{titlepage}
  \thispagestyle{empty}
  \newpagecolor{crema}
  \setlength{\parindent}{0pt}
  \begin{tikzpicture}[remember picture,overlay]
    % Brillo de sol como en el fondo de la web (detrás de la cabecera)
    \shade[inner color=limon!70, outer color=crema] ([xshift=-3.5cm,yshift=-9cm]current page.north east) circle (4cm);
    % Cabecera verde de hoja con la franja de toldo
    \fill[hoja] (current page.north west) rectangle ([yshift=-5.4cm]current page.north east);
    % Paisaje de la huerta al pie de la portada
    \node[anchor=south, inner sep=0, opacity=0.8] at (current page.south)
      {\includegraphics[width=\paperwidth]{img/huerta.pdf}};
  \end{tikzpicture}
  \toldo{5.4cm}

  \vspace*{-0.4cm}
  {\color{white}
    \large\textsc{Universidad Nacional de Educación a Distancia}\\[0.4cm]
    {\titular\Large \asignatura}\par}

  \vspace{3.6cm}
  {\color{naranja}\bfseries\large COOPERATIVA HUERTA DE MURCIA · BASES DE DATOS\par}
  \vspace{0.2cm}
  {\titular\fontsize{58}{60}\selectfont\color{tierra} Misión \textcolor{naranja}{SQL}\par}
  \vspace{0.3cm}
  {\Large\color{hoja} El enigma de la Cooperativa Huerta de Murcia\par}
  \vspace{0.4cm}
  {\large Recurso educativo gamificado: presentación interactiva\\
   y tres juegos digitales sobre el lenguaje SQL\par}

  \vspace{0.9cm}
  \begin{tcolorbox}[colback=papel, colframe=borde, arc=4mm, boxrule=0.6pt, drop fuzzy shadow=borde]
    \begin{tabularx}{\linewidth}{@{}>{\leavevmode\bfseries\color{hoja}}p{4.2cm}X@{}}
      Actividad        & \actividad \\
      Bloque temático  & \bloque \\
      Materia de FP    & \materiaFP \\
      Contextualización & Comunidad Autónoma de la Región de Murcia \\
    \end{tabularx}
  \end{tcolorbox}

  \vspace{0.5cm}
  \begin{tabularx}{\linewidth}{@{}>{\leavevmode\bfseries\color{hoja}}p{4.2cm}X@{}}
    Estudiante      & \alumno \\
    DNI             & \dni \\
    Correo electrónico & \href{mailto:\email}{\email} \\
    Centro asociado & \centro \\
    Curso académico & \curso \\
  \end{tabularx}
  \vfill
\end{titlepage}
\restorepagecolor

\tableofcontents
\newpage

% =====================================================================
\section{Planteamiento general}
% =====================================================================

\subsection{Justificación de la temática}

El lenguaje SQL (\emph{Structured Query Language}) es el lenguaje estándar para definir,
consultar y manipular la información almacenada en sistemas gestores de bases de datos
relacionales. Se basa en el modelo relacional propuesto por Codd \cite{codd1970}, nació en
IBM con el nombre de SEQUEL \cite{chamberlin1974} y está normalizado internacionalmente en
la serie ISO/IEC~9075 \cite{iso9075}. Es el lenguaje de los SGBD más usados en la industria,
como MySQL \cite{mysqldoc}, PostgreSQL \cite{postgresdoc} o SQLite \cite{sqlitedoc}, y la
bibliografía de referencia de la materia lo trata como contenido central
\cite{silberschatz,elmasri}. Prácticamente cualquier aplicación que desarrolle o administre un técnico
superior de la familia de Informática y Comunicaciones (aplicaciones web, móviles, de
escritorio o servicios en red) acaba leyendo y escribiendo datos mediante SQL, por lo que
su dominio es una competencia transversal y muy demandada por el tejido empresarial. La
encuesta anual de Stack Overflow lo sitúa de forma constante entre los lenguajes más
utilizados por los profesionales del desarrollo de software \cite{stackoverflow2025}.

Sin embargo, la investigación sobre la enseñanza de SQL muestra que su
aprendizaje presenta dificultades recurrentes \cite{taipalus2020}. Los
estudios que analizan las consultas escritas por estudiantes identifican
errores sistemáticos tanto sintácticos como semánticos y lógicos
\cite{taipalus2018}. Los más frecuentes aparecen en el agrupamiento
(\texttt{GROUP BY}/\texttt{HAVING}), en las composiciones de varias
tablas, en las composiciones externas y en las subconsultas
\cite{ahadi2016,taipalus2018}. Influyen también la naturaleza declarativa
del lenguaje, poco intuitiva para un alumnado que empieza a programar en
lenguajes imperativos, y el hecho de que el orden lógico de evaluación
de las cláusulas no coincide con el orden en que se escriben
\cite{taipalus2018}.

Estas dificultades ponen de manifiesto que aprender SQL no consiste
únicamente en memorizar su sintaxis, sino también en desarrollar la
capacidad de traducir problemas de consulta de datos a expresiones
correctas. Para ello, resulta necesario que el alumnado se enfrente
progresivamente a problemas de distinta complejidad, formule consultas,
compruebe sus resultados e identifique y corrija sus errores. En este
proceso, la retroalimentación desempeña un papel fundamental, ya que
permite al estudiante reconocer sus errores y orientar sus siguientes
intentos \cite{hattie2007}. De hecho, una de las líneas de investigación
en la enseñanza de SQL se centra en el desarrollo de herramientas
capaces de evaluar automáticamente las consultas y proporcionar
retroalimentación inmediata \cite{taipalus2020}.

No obstante, disponer de oportunidades para practicar no garantiza
que el alumnado las aproveche de forma continuada. La resolución
reiterada de consultas, especialmente cuando estas presentan una
dificultad creciente o requieren varios intentos, exige mantener
la implicación y el esfuerzo durante el proceso de aprendizaje.
Por tanto, además de proporcionar ejercicios y retroalimentación,
resulta conveniente adoptar estrategias didácticas que favorezcan
la participación activa y la persistencia del alumnado ante las
dificultades.

En este contexto, la gamificación constituye una posible respuesta
a esta necesidad. Entendida como la incorporación de elementos de
diseño de juegos en contextos no lúdicos \cite{deterding2011},
permite organizar la práctica mediante retos progresivos, objetivos
claros, sistemas de puntuación y mecanismos de progresión que hacen
visible el avance del alumnado \cite{kapp2012,werbach2012}.
Aplicados a la enseñanza de SQL, estos elementos permiten estructurar
las actividades en niveles de dificultad creciente, reconocer los
logros alcanzados y estimular nuevos intentos tras cometer errores.
De este modo, la gamificación podríar proporcionar un marco para favorecer
la participación y la continuidad de la práctica y retroalimentación en el contexto
de la enseñanza de SQL.
La investigación previa sobre gamificación educativa señala, en
general, efectos positivos sobre la motivación, la participación
y los resultados de aprendizaje \cite{hamari2014,dicheva2015,sailer2020}. 

En resumen, las dificultades de aprendizaje de SQL justifican
la necesidad de una práctica progresiva, continuada y acompañada
de retroalimentación. A su vez, la gamificación ofrece una estrategia
para fomentar la implicación del alumnado en ese proceso mediante
retos, objetivos y mecanismos de progresión. Sobre esta base, se han
seleccionado las \emph{sentencias de consulta y de manipulación de datos}
de SQL como eje de un recurso educativo gamificado.

\noindent\textbf{Valor añadido.} Además, para que la práctica tenga sentido para el alumnado, todos los retos se plantean
sobre un caso cercano: una cooperativa agrícola ficticia de la Región de Murcia, con socios
de Murcia, Lorca, Cartagena o Cieza y productos como limones, alcachofas o melocotones.
Esta decisión es coherente con el currículo autonómico vigente, que establece que el
currículo se implantará «teniendo en cuenta la realidad socioeconómica, las necesidades del
tejido empresarial y las características geográficas, socio-productivas y laborales propias
del entorno» y reconoce a los centros autonomía pedagógica para adaptar las enseñanzas «a las
características concretas del entorno socioeconómico, cultural y profesional»
\cite[art.~4]{murcia2026}.

\subsection{Destinatarios y módulo profesional}

El recurso está dirigido al alumnado de \textbf{primer curso} de los Ciclos Formativos de
Grado Superior (CFGS) de \textbf{Desarrollo de Aplicaciones Multiplataforma (DAM)} y
\textbf{Desarrollo de Aplicaciones Web (DAW)}, de la familia profesional Informática y
Comunicaciones, que cursa el módulo profesional \textbf{0484 Bases de datos} en un centro de
la Región de Murcia. Este módulo es común a ambos títulos \cite{rd450,rd686,rd405,rd500}.
Sus elementos curriculares se fijan en el RD 405/2023 \cite{rd405} y su organización en la
región, en la Orden de 21 de septiembre de 2026 \cite{murcia2026}. La
\autoref{tab:modulo} resume sus datos básicos.

\begin{table}[H]
\centering
\begin{tabularx}{\linewidth}{@{}>{\bfseries}p{0.3\linewidth}X@{}}
\toprule
Código y denominación & 0484 Bases de datos \\
Ciclos formativos & CFGS Desarrollo de Aplicaciones Multiplataforma (DAM) y CFGS
                    Desarrollo de Aplicaciones Web (DAW) \\
Curso & Primero \cite[anexos XLIV y XLV]{murcia2026} \\
Duración en la Región de Murcia & 175 horas, 5 horas semanales
                    \cite[anexos XLIV y XLV]{murcia2026} \\
Duración mínima estatal & 105 horas \cite{rd405} \\
Créditos ECTS & 12 \cite{rd405} \\
Resultados de aprendizaje & 7 (los trabajados, en la \autoref{tab:ra}) \\
\bottomrule
\end{tabularx}
\caption{Datos básicos del módulo profesional 0484}
\label{tab:modulo}
\end{table}

\paragraph{Perfil del alumnado.} Se trata de alumnado mayor de edad o próximo a serlo, con
perfiles de acceso heterogéneos (Bachillerato, Grado Medio, prueba de acceso o
reincorporación al sistema educativo, según las vías de acceso a los ciclos de grado
superior que regula la normativa de Formación Profesional \cite{lofp,rd659}), con
conocimientos básicos de ofimática y que en paralelo inicia el módulo de Programación.

\paragraph{Función del módulo y perfil profesional.} Según las orientaciones pedagógicas del
RD 405/2023, el módulo \textit{contiene la formación necesaria para desempeñar la función de
programación de bases de datos}, que incluye la planificación y realización del diseño
lógico y físico de una base de datos, la normalización de esquemas, la inserción y
manipulación de datos, la planificación y realización de consultas y la programación de
procedimientos almacenados. Las actividades profesionales asociadas se aplican en la gestión
de la información almacenada en bases de datos y en el desarrollo de aplicaciones que acceden
a ellas \cite{rd405}.

En el título de DAM, el módulo contribuye a alcanzar los objetivos generales c), e), f), p)
y t) del ciclo y las competencias b), c), e), p) y t) del título. Entre ellas destacan la
competencia c), \textit{Gestionar bases de datos, interpretando su diseño lógico y verificando
integridad, consistencia, seguridad y accesibilidad de los datos}, y la e), \textit{Desarrollar
aplicaciones multiplataforma con acceso a bases de datos utilizando lenguajes, librerías y
herramientas adecuados a las especificaciones}. En el título de DAW contribuye a los
objetivos generales e), f), g), p) y r) y a las competencias d), e) y p) \cite{rd405}.

\paragraph{Estructura del módulo.} El módulo se organiza en siete resultados de aprendizaje,
que recorren el ciclo completo de trabajo con bases de datos: reconocer los sistemas de
almacenamiento y los SGBD (RA1), crear bases de datos relacionales (RA2), consultar (RA3)
y modificar (RA4) la información, programar procedimientos almacenados (RA5), diseñar
modelos relacionales normalizados a partir de diagramas entidad/relación (RA6) y gestionar
bases de datos no relacionales (RA7). Sus contenidos básicos se agrupan en los bloques
\emph{Almacenamiento de la información}, \emph{Bases de datos relacionales},
\emph{Realización de consultas}, \emph{Tratamiento de datos}, \emph{Programación de bases de
datos}, \emph{Interpretación de diagramas entidad/relación} y \emph{Uso de bases de datos no
relacionales} \cite{rd405}.

\paragraph{Lugar de la actividad en el módulo.} El recurso se sitúa en el centro de esa
secuencia. Se presupone que el alumnado ya ha trabajado el RA2, \textit{Crea bases de datos
definiendo su estructura y las características de sus elementos según el modelo
relacional}, es decir, que conoce el modelo relacional y sabe crear tablas con sus claves y
restricciones. Con este recurso se inicia en la consulta y la modificación de datos con SQL
(RA3 y RA4, \autoref{sec:ra}), antes de pasar a la programación de procedimientos
almacenados (RA5). %Responde así a una de las líneas de actuación que las orientaciones
%pedagógicas proponen para el módulo: «la realización de operaciones de consulta y
%modificación sobre los datos almacenados» \cite{rd405}.

\paragraph{Uso en otros módulos.} El recurso también se puede usar, con cambios menores, en
el módulo \emph{0372 Gestión de bases de datos} (CFGS ASIR), también de primer curso
\cite{rd1629,rd500,ordenASIR}\cite[anexo XLIII]{murcia2026}, cuyos contenidos de consulta y
manipulación de datos son equivalentes, y como repaso al inicio del módulo
\emph{0486 Acceso a datos} (segundo curso del CFGS DAM) \cite{rd405}\cite[anexo XLIV]{murcia2026}.

\subsection{Objetivos}

\paragraph{Objetivo general.} Que el alumnado sea capaz de consultar y modificar la
información de una base de datos relacional mediante sentencias SQL correctas y eficientes,
valorando distintas alternativas para resolver un mismo problema.

\paragraph{Objetivos específicos.} Al finalizar la actividad, el alumnado será capaz de:
\begin{enumerate}[label=\textbf{O\arabic*.}]
  \item Distinguir los sublenguajes de SQL (DDL, DML, DCL y control de transacciones) y
        clasificar sus sentencias principales.
  \item Construir consultas simples sobre una tabla con selección (\texttt{WHERE}),
        proyección, ordenación (\texttt{ORDER BY}) y eliminación de duplicados
        (\texttt{DISTINCT}), usando operadores relacionales, lógicos, \texttt{BETWEEN},
        \texttt{IN}, \texttt{LIKE} e \texttt{IS NULL}.
  \item Obtener valores de resumen con funciones de agregado, agrupando
        (\texttt{GROUP BY}) y filtrando grupos (\texttt{HAVING}).
  \item Combinar información de varias tablas mediante composiciones internas
        (\texttt{INNER JOIN}) y externas (\texttt{LEFT}/\texttt{RIGHT JOIN}).
  \item Resolver consultas mediante subconsultas y compararlas con soluciones equivalentes
        basadas en composiciones.
  \item Insertar, actualizar y borrar datos (\texttt{INSERT}, \texttt{UPDATE},
        \texttt{DELETE}) respetando la integridad referencial.
  \item Reconocer el funcionamiento de las transacciones (\texttt{START TRANSACTION},
        \texttt{COMMIT}, \texttt{ROLLBACK}, \texttt{SAVEPOINT}).
  \item Mostrar autonomía, constancia y espíritu de superación ante los errores,
        y colaborar con el grupo durante los retos.
\end{enumerate}

\subsection{Marco normativo de aplicación}

\subsubsection{Normativa estatal}

\begin{itemize}
  \item \textbf{Ley Orgánica 2/2006, de 3 de mayo, de Educación} (LOE), modificada por la
        Ley Orgánica 3/2020, de 29 de diciembre (LOMLOE) \cite{loe,lomloe}.
  \item \textbf{Ley Orgánica 3/2022, de 31 de marzo, de ordenación e integración de la
        Formación Profesional} \cite{lofp}, y su desarrollo mediante el
        \textbf{Real Decreto 659/2023, de 18 de julio}, por el que se desarrolla la
        ordenación del Sistema de Formación Profesional \cite{rd659}.
  \item \textbf{Real Decreto 278/2023, de 11 de abril}, por el que se establece el calendario
        de implantación del Sistema de Formación Profesional establecido por la Ley Orgánica
        3/2022 \cite{rd278}.
  \item \textbf{Real Decreto 450/2010, de 16 de abril}, por el que se establece el título
        de Técnico Superior en Desarrollo de Aplicaciones Multiplataforma y se fijan sus
        enseñanzas mínimas \cite{rd450}, y \textbf{Real Decreto 686/2010, de 20 de
        mayo}, para el título de Técnico Superior en Desarrollo de Aplicaciones Web
        \cite{rd686}. Ambos siguen vigentes con las dos modificaciones siguientes.
  \item \textbf{Real Decreto 405/2023, de 29 de mayo}, por el que se actualizan los títulos
        de Técnico Superior en Desarrollo de Aplicaciones Multiplataforma y Técnico Superior
        en Desarrollo de Aplicaciones Web y se fijan sus enseñanzas mínimas. Da nueva
        redacción al módulo 0484 y establece sus \emph{resultados de aprendizaje, criterios
        de evaluación y contenidos básicos} \cite{rd405}.
  \item \textbf{Real Decreto 500/2024, de 21 de mayo}, por el que se modifican determinados
        reales decretos por los que se establecen títulos de Formación Profesional de grado
        superior y se fijan sus enseñanzas mínimas. Adapta los títulos de DAM, DAW y ASIR,
        entre otros, a la nueva ordenación del RD 659/2023 \cite{rd500}.
  \item \textbf{Orden EDU/2000/2010, de 13 de julio}, por la que se establece el currículo
        del ciclo formativo de grado superior de DAM \cite{ordenDAM}, y
        \textbf{Orden EDU/2887/2010, de 2 de noviembre}, para el ciclo de DAW
        \cite{ordenDAW}. Desarrollan el currículo en el ámbito de gestión del Ministerio,
        por lo que no se aplican en la Región de Murcia; se citan solo como referencia.
\end{itemize}

\subsubsection{Normativa autonómica (Región de Murcia)}

\begin{itemize}
  \item \textbf{Decreto n.º 13/2026, de 5 de marzo}, por el que se establece la ordenación
        de la Formación Profesional de Grado D y de Grado E en la Comunidad Autónoma de la
        Región de Murcia. Determina la estructura mínima de los currículos de los ciclos
        formativos \cite{decreto13}.
  \item \textbf{Orden de 21 de septiembre de 2026}, de la Consejería de Educación y Formación
        Profesional, por la que se establecen los currículos de los títulos de los ciclos
        formativos de formación profesional de grado superior y se regulan determinados
        aspectos organizativos de aquellos en la Comunidad Autónoma de la Región de Murcia
        (BORM n.º~222, de 25 de septiembre de 2026). Produce efectos académicos desde el
        curso 2024/2025 y deroga las Órdenes de 12 de marzo de 2013 que establecían los
        currículos de DAM y DAW en la región \cite{murcia2026}.
\end{itemize}

La Orden de 21 de septiembre de 2026 establece que los resultados de aprendizaje y los
criterios de evaluación de cada módulo son los del Anexo I de los reales decretos de
título y tienen carácter \emph{prescriptivo}, mientras que los contenidos básicos tienen
carácter \emph{orientativo} y corresponde a los equipos docentes concretarlos en las
programaciones \cite[art.~3]{murcia2026}. Además, dispone que el currículo se adapte a la
realidad socioeconómica y productiva del entorno \cite[art.~4]{murcia2026}. Sus anexos
XLIV (DAM) y XLV (DAW) sitúan el módulo 0484 en el primer curso, con 175 horas y
5 horas semanales. Por ese motivo, el recurso toma como referencia los resultados de
aprendizaje y criterios de evaluación del RD 405/2023 y usa como caso práctico una
cooperativa hortofrutícola murciana, un sector económico clave en la región.

\subsubsection{Resultados de aprendizaje y criterios de evaluación}
\label{sec:ra}

El módulo 0484 \emph{Bases de datos} tiene siete resultados de aprendizaje (RA), con sus
criterios de evaluación (CE), en la redacción dada por el Anexo~I del RD 405/2023
\cite{rd405}. En la Región de Murcia son prescriptivos para el diseño de las programaciones y
la evaluación \cite[art.~3]{murcia2026}. La \autoref{tab:ra} recoge los tres que
la actividad trabaja o repasa (RA2, RA3 y RA4) y marca con \trab{} los criterios concretos
que se abordan. Los otros cuatro (RA1, RA5, RA6 y RA7) no se trabajan.

\begin{longtable}{@{}p{0.27\linewidth}p{0.68\linewidth}@{}}
\caption{Resultados de aprendizaje y criterios de evaluación del módulo 0484 (RD 405/2023) que se trabajan o repasan}
\label{tab:ra}\\
\toprule
\textbf{Resultado de aprendizaje} & \textbf{Criterios de evaluación} \\
\midrule
\endfirsthead
\toprule
\textbf{Resultado de aprendizaje} & \textbf{Criterios de evaluación} \\
\midrule
\endhead
\bottomrule
\endfoot
\textbf{RA2.} Crea bases de datos definiendo su estructura y las características de sus
elementos según el modelo relacional.

\smallskip\emph{Se repasa}: claves y restricciones, necesarias para entender los errores de
integridad referencial de los retos.
&
\begin{enumerate}[label=\alph*), leftmargin=*, nosep]
  \item Se ha analizado el formato de almacenamiento de la información.
  \item Se han creado las tablas y las relaciones entre ellas.
  \item Se han seleccionado los tipos de datos adecuados.
  \item Se han definido los campos clave en las tablas. \trab
  \item Se han implantado las restricciones reflejadas en el diseño lógico. \trab
  \item Se han creado vistas.
  \item Se han creado los usuarios y se les han asignado privilegios.
  \item Se han utilizado asistentes, herramientas gráficas y los lenguajes de definición
        y control de datos.
\end{enumerate} \\
\textbf{RA3.} Consulta la información almacenada en una base de datos empleando
asistentes, herramientas gráficas y el lenguaje de manipulación de datos.

\smallskip\emph{Se trabaja}: núcleo de la actividad.
&
\begin{enumerate}[label=\alph*), leftmargin=*, nosep]
  \item Se han identificado las herramientas y sentencias para realizar consultas. \trab
  \item Se han realizado consultas simples sobre una tabla. \trab
  \item Se han realizado consultas sobre el contenido de varias tablas mediante
        composiciones internas. \trab
  \item Se han realizado consultas sobre el contenido de varias tablas mediante
        composiciones externas. \trab
  \item Se han realizado consultas resumen. \trab
  \item Se han realizado consultas con subconsultas. \trab
  \item Se han realizado consultas que implican múltiples selecciones.
  \item Se han aplicado criterios de optimización de consultas. \trab
\end{enumerate} \\
\textbf{RA4.} Modifica la información almacenada en la base de datos utilizando
asistentes, herramientas gráficas y el lenguaje de manipulación de datos.

\smallskip\emph{Se trabaja}: núcleo de la actividad.
&
\begin{enumerate}[label=\alph*), leftmargin=*, nosep]
  \item Se han identificado las herramientas y sentencias para modificar el contenido
        de la base de datos. \trab
  \item Se han insertado, borrado y actualizado datos en las tablas. \trab
  \item Se ha incluido en una tabla la información resultante de la ejecución de una
        consulta. \trab
  \item Se han diseñado guiones de sentencias para llevar a cabo tareas complejas. \trab
  \item Se ha reconocido el funcionamiento de las transacciones. \trab
  \item Se han anulado parcial o totalmente los cambios producidos por una transacción.
        \trab
  \item Se han identificado los efectos de las distintas políticas de bloqueo de
        registros. \trab
  \item Se han adoptado medidas para mantener la integridad y consistencia de la
        información. \trab
\end{enumerate} \\
\end{longtable}

En resumen, la actividad se centra en los \textbf{RA3} y \textbf{RA4}. Trabaja todos sus
criterios excepto el 3.g, y repasa los criterios 2.d y 2.e del \textbf{RA2}.

\subsection{Contenidos que se trabajan}

De los contenidos básicos del módulo \cite{rd405}, que en la Región de Murcia tienen carácter
orientativo \cite[art.~3]{murcia2026}, el recurso trabaja los bloques \emph{Realización de
consultas} y \emph{Tratamiento de datos}, además de los sublenguajes DDL y DCL del bloque
\emph{Bases de datos relacionales}:

\begin{table}[H]
\centering
\begin{tabularx}{\linewidth}{@{}p{0.3\linewidth}X@{}}
\toprule
\textbf{Bloque de contenidos} & \textbf{Contenidos concretos} \\
\midrule
Introducción al lenguaje SQL &
Sublenguajes DDL, DML, DCL y de control de transacciones. Herramientas gráficas y de línea
de comandos para ejecutar sentencias. \\
\addlinespace
Realización de consultas &
Proyección, selección y ordenación de registros (\texttt{SELECT}, \texttt{WHERE},
\texttt{ORDER BY}, alias y \texttt{DISTINCT}). Operadores de comparación y lógicos,
\texttt{BETWEEN}, \texttt{IN}, \texttt{LIKE}, \texttt{IS NULL}. Consultas de resumen:
funciones de agregado (\texttt{COUNT}, \texttt{SUM}, \texttt{AVG}, \texttt{MIN},
\texttt{MAX}). Agrupamiento de registros (\texttt{GROUP BY}) y filtrado de grupos
(\texttt{HAVING}). Composiciones internas y externas. Subconsultas. Optimización de
consultas: comparación de soluciones equivalentes. \\
\addlinespace
Tratamiento de datos &
Inserción, borrado y modificación de registros. Integridad referencial. Subconsultas y
composiciones en órdenes de edición. Transacciones: \texttt{COMMIT}, \texttt{ROLLBACK},
\texttt{SAVEPOINT}. Políticas de bloqueo y concurrencia (nociones). \\
\addlinespace
Contenidos transversales &
Resolución de problemas, pensamiento computacional, trabajo en equipo, tolerancia al error
y conocimiento del entorno productivo de la Región de Murcia. \\
\bottomrule
\end{tabularx}
\caption{Contenidos trabajados en la actividad}
\end{table}

% =====================================================================
\section{Descripción de la actividad}
% =====================================================================

\subsection{Visión general y narrativa}

\emph{Misión SQL: el enigma de la Cooperativa Huerta de Murcia} es un recurso gamificado
formado por:

\begin{enumerate}
  \item Una \textbf{presentación interactiva} en Genially, que sirve de hilo conductor
        y donde se desarrollan la narrativa, los contenidos teóricos y la sala final.
  \item \textbf{Tres juegos digitales educativos}, desarrollados como páginas web propias
        (HTML, CSS y JavaScript) y publicados con GitHub Pages. La presentación los enlaza
        en el orden de la misión.
\end{enumerate}

\paragraph{Narrativa.} Una tormenta de granizo ha dañado el servidor de la
\emph{Cooperativa Huerta de Murcia} justo antes de la campaña de exportación. La gerencia
ha contratado al alumnado como \emph{técnicos junior de bases de datos} para recuperar la
información de socios, productos y entregas. El acceso a la copia de seguridad está protegido
por una \textbf{caja fuerte digital con un código de cuatro dígitos}. Cada juego superado
revela uno o dos dígitos (la «llave» del juego), y solo el equipo que reúna los cuatro
podrá abrir la caja fuerte en la última diapositiva y «salvar la campaña».

\paragraph{Elementos de gamificación.} Siguiendo el marco de mecánicas, dinámicas y
componentes de Werbach y Hunter \cite{werbach2012}:

\begin{itemize}
  \item \textbf{Dinámicas}: narrativa, progresión por niveles (de lo sencillo a lo complejo),
        sensación de logro y colaboración por equipos.
  \item \textbf{Mecánicas}: retos, retroalimentación inmediata, recompensas (dígitos
        de la llave), competición sana entre equipos y límite de tiempo en el juego 2.
  \item \textbf{Componentes}: puntos, insignias por juego («Clasificador»,
        «Detective de consultas», «Maestro del DML»), tabla de clasificación y desbloqueo
        de la caja fuerte final.
\end{itemize}

\begin{figure}[H]
\centering
\begin{tikzpicture}[
  node distance=0.9cm and 0.7cm,
  caja/.style={rectangle, rounded corners=3pt, draw=principal, very thick,
               fill=principal!8, minimum width=2.6cm, minimum height=1.25cm,
               align=center, font=\small},
  llave/.style={rectangle, rounded corners=3pt, draw=acento, thick, fill=acento!10,
               font=\small\bfseries, align=center},
  flecha/.style={-{Stealth[length=3mm]}, thick, principal}
]
  \node[caja] (g0) {Genially\\Narrativa y teoría};
  \node[caja, right=of g0] (j1) {Juego 1\\Clasifica\\y ordena};
  \node[caja, right=of j1] (j2) {Juego 2\\¿Qué devuelve\\la consulta?};
  \node[caja, right=of j2] (j3) {Juego 3\\Terminal de\\la Cooperativa};
  \node[caja, below=1.6cm of j2, draw=acento, fill=acento!10] (fin) {Genially\\Caja fuerte final};
  \draw[flecha] (g0) -- (j1);
  \draw[flecha] (j1) -- (j2);
  \draw[flecha] (j2) -- (j3);
  \node[llave, below=0.25cm of j1] (k1) {dígito 1};
  \node[llave, below=0.25cm of j2] (k2) {dígito 2};
  \node[llave, below=0.25cm of j3] (k3) {dígitos 3 y 4};
  \draw[flecha, acento] (k1) |- (fin);
  \draw[flecha, acento] (k2) -- (fin);
  \draw[flecha, acento] (k3) |- (fin);
\end{tikzpicture}
\caption{Itinerario de la misión y obtención del código de la caja fuerte}
\end{figure}

\subsection{Base de datos de trabajo}

Todos los retos se resuelven sobre la misma base de datos, \texttt{cooperativa}, que el
alumnado conoce en la presentación y que el juego~3 carga en el navegador. Tiene tres
tablas:

\begin{itemize}
  \item \texttt{socio(\underline{id\_socio}, nombre, municipio, fecha\_alta)}
  \item \texttt{producto(\underline{id\_producto}, nombre, categoria, precio\_kg)}
  \item \texttt{entrega(\underline{id\_entrega}, id\_socio\textsuperscript{FK},
        id\_producto\textsuperscript{FK}, fecha, kilos)}
\end{itemize}

\begin{table}[H]
\centering\small
\begin{tabular}{@{}clll@{}}
\toprule
\multicolumn{4}{c}{\texttt{socio}} \\
\midrule
\textbf{id} & \textbf{nombre} & \textbf{municipio} & \textbf{fecha\_alta} \\
1 & Ana Martínez   & Murcia    & 2019-02-11 \\
2 & José López     & Lorca     & 2020-05-03 \\
3 & Carmen Sánchez & Cartagena & 2021-09-20 \\
4 & Pedro Ruiz     & Murcia    & 2022-01-15 \\
5 & Lucía Gómez    & Cieza     & 2025-11-30 \\
\bottomrule
\end{tabular}
\hfill
\begin{tabular}{@{}clll@{}}
\toprule
\multicolumn{4}{c}{\texttt{producto}} \\
\midrule
\textbf{id} & \textbf{nombre} & \textbf{categoria} & \textbf{precio\_kg} \\
1 & Limón       & Cítrico        & 0,60 \\
2 & Naranja     & Cítrico        & 0,45 \\
3 & Pimiento    & Hortaliza      & 1,20 \\
4 & Alcachofa   & Hortaliza      & 1,80 \\
5 & Melocotón   & Fruta de hueso & 1,50 \\
6 & Albaricoque & Fruta de hueso & 2,10 \\
\bottomrule
\end{tabular}

\vspace{0.4cm}
\begin{tabular}{@{}ccccr@{}}
\toprule
\multicolumn{5}{c}{\texttt{entrega}} \\
\midrule
\textbf{id\_entrega} & \textbf{id\_socio} & \textbf{id\_producto} & \textbf{fecha} & \textbf{kilos} \\
1 & 1 & 1 & 2026-03-02 & 500 \\
2 & 1 & 3 & 2026-03-05 & 200 \\
3 & 2 & 2 & 2026-03-05 & 800 \\
4 & 2 & 4 & 2026-03-10 & 150 \\
5 & 3 & 1 & 2026-03-12 & 300 \\
6 & 3 & 5 & 2026-06-20 & 400 \\
7 & 4 & 3 & 2026-03-15 & 250 \\
8 & 4 & 4 & 2026-03-18 & 100 \\
9 & 1 & 5 & 2026-06-25 & 350 \\
\bottomrule
\end{tabular}
\caption{Contenido de la base de datos \texttt{cooperativa}}
\label{tab:datos}
\end{table}

Los datos están pensados para que aparezcan casos didácticos concretos: una socia sin
entregas (Lucía Gómez), útil para las composiciones externas; un producto sin entregas
(Albaricoque), útil para \texttt{NOT IN}; y dos socios del mismo municipio, útil para
\texttt{DISTINCT}. El guion completo de creación se incluye en el
\autoref{ap:script}.

\subsection{Presentación interactiva (Genially)}

\begin{nota}
\textbf{URL pública:} \url{\urlGenially}\\
\textbf{Número de diapositivas:} \numDiapositivas
\end{nota}

La presentación es el tablero de la misión. Combina diapositivas de narrativa,
diapositivas de contenido con ejemplos interactivos (ventanas emergentes con el resultado de
cada consulta y botones «¿Y si…?») y botones de acceso a cada juego. Su estructura de
contenidos es la siguiente:

\begin{longtable}{@{}p{0.2\linewidth}p{0.75\linewidth}@{}}
\toprule
\textbf{Sección} & \textbf{Contenido} \\
\midrule
\endhead
Portada y misión & Título, presentación de la Cooperativa Huerta de Murcia, la tormenta
y el encargo al equipo técnico. Normas del juego, formación de equipos y explicación de la
caja fuerte. \\
\addlinespace
Objetivos & Qué vamos a aprender: objetivos O1–O8 en lenguaje cercano al alumnado. \\
\addlinespace
La base de datos & Esquema relacional de \texttt{cooperativa} con claves primarias y
ajenas; ventanas emergentes con el contenido de cada tabla. \\
\addlinespace
¿Qué es SQL? & Historia breve, estándar ISO/IEC 9075, SGBD habituales (MySQL/MariaDB,
PostgreSQL, Oracle, SQLite). Sublenguajes DDL, DML, DCL y TCL con ejemplos. \\
\addlinespace
\textbf{Acceso al juego 1} & Botón «Clasifica y ordena» y casilla para anotar el
\emph{dígito~1}. \\
\addlinespace
Consultas simples & Estructura de \texttt{SELECT}; proyección, alias, \texttt{DISTINCT};
\texttt{WHERE} con operadores de comparación y lógicos; \texttt{BETWEEN}, \texttt{IN},
\texttt{LIKE}, \texttt{IS NULL}; \texttt{ORDER BY}. Orden de escritura frente a orden lógico
de evaluación de las cláusulas. \\
\addlinespace
Consultas de resumen & Funciones de agregado; \texttt{GROUP BY}; diferencia entre
\texttt{WHERE} y \texttt{HAVING}, con una animación de cómo se forman los grupos. \\
\addlinespace
Varias tablas & Producto cartesiano; \texttt{INNER JOIN}; \texttt{LEFT}/\texttt{RIGHT JOIN}
con diagramas de conjuntos; ejemplo de la socia sin entregas. \\
\addlinespace
Subconsultas & Subconsultas escalares, con \texttt{IN}/\texttt{NOT IN} y
\texttt{EXISTS}; comparación entre una subconsulta y un \texttt{JOIN} equivalente
(CE~3.h). \\
\addlinespace
\textbf{Acceso al juego 2} & Botón «¿Qué devuelve la consulta?» y casilla del
\emph{dígito~2}. \\
\addlinespace
Manipulación de datos & \texttt{INSERT} (incluido \texttt{INSERT \ldots{} SELECT}),
\texttt{UPDATE}, \texttt{DELETE}; peligro de olvidar el \texttt{WHERE}; errores de
integridad referencial y opciones \texttt{ON DELETE}/\texttt{ON UPDATE}. \\
\addlinespace
Transacciones & Propiedades ACID; \texttt{START TRANSACTION}, \texttt{COMMIT},
\texttt{ROLLBACK}, \texttt{SAVEPOINT}; nociones de bloqueo y acceso concurrente. \\
\addlinespace
\textbf{Acceso al juego 3} & Botón «Terminal de la Cooperativa» y casillas de los
\emph{dígitos~3 y~4}. \\
\addlinespace
Caja fuerte final & Teclado numérico interactivo: el código correcto muestra la pantalla
de victoria y la insignia final; un código erróneo devuelve una pista sobre el juego que hay
que revisar. \\
\addlinespace
Cierre & Resumen de lo aprendido, autoevaluación, enlace a la base de datos para
practicar en casa, créditos y bibliografía. \\
\bottomrule
\end{longtable}

\subsection{Juegos digitales educativos}

\begin{nota}
Los tres juegos se publican juntos, en una única aplicación de GitHub Pages:\\
\textbf{Repositorio:} \url{\urlRepo}\\
\textbf{Página índice de los juegos:} \url{\urlJuegos}
\end{nota}

Los juegos forman una \textbf{aplicación web común} con \textbf{tres juegos diferenciados}.
La aplicación ofrece la interfaz, la narrativa y las mecánicas de gamificación
compartidas: perfil del equipo, marcador global de puntos, insignias, mapa del itinerario y
registro de los dígitos de la llave. Cada juego es un juego independiente, con su propia
naturaleza, mecánica y forma de interacción. No son fases de un mismo juego. El
itinerario es progresivo: cada juego trabaja un nivel de complejidad mayor que el anterior
y se desbloquea al superar el previo, aunque todos se pueden volver a jugar por separado.

\begin{table}[H]
\centering\small
\begin{tabularx}{\linewidth}{@{}l>{\raggedright}p{2.7cm}>{\raggedright}p{3.2cm}>{\raggedright\arraybackslash}X@{}}
\toprule
\textbf{Juego} & \textbf{Tipo de juego} & \textbf{Interacción} & \textbf{Habilidad que se trabaja} \\
\midrule
1. Clasifica y ordena & Puzle de clasificación y ordenación & Arrastrar y soltar piezas &
Reconocer la sintaxis y la estructura del lenguaje (\emph{recordar y comprender}) \\
\addlinespace
2. ¿Qué devuelve la consulta? & Concurso de preguntas contrarreloj & Elegir una respuesta con
cuenta atrás & Leer e interpretar sentencias y predecir su resultado (\emph{comprender y
analizar}) \\
\addlinespace
3. Terminal de la Cooperativa & Simulación o \emph{escape room} de resolución de
incidencias & Escribir y ejecutar SQL real en una consola & Escribir sentencias que
resuelven problemas reales (\emph{aplicar y crear}) \\
\bottomrule
\end{tabularx}
\caption{Naturaleza diferenciada de los tres juegos}
\end{table}

La aplicación está formada por páginas web estáticas (HTML5, CSS y JavaScript sin servidor),
así que se puede publicar gratis en GitHub Pages, funciona en cualquier navegador de
ordenador, tableta o móvil, y no pide registro ni recoge datos personales del alumnado (puntuaciones y
progreso se guardan solo en el navegador con \texttt{localStorage}). Todos los juegos
comparten la estética de la presentación y el mismo esquema:
pantalla de inicio con la parte de la narrativa que corresponde, instrucciones, desarrollo
con retroalimentación inmediata y pantalla final con puntuación, insignia y dígito de la llave.

% ---------------------------------------------------------------------
\subsubsection{Juego 1: «Clasifica y ordena»}
% ---------------------------------------------------------------------

\begin{nota}
\textbf{URL pública:} \url{\urlJuegoUno}\\
\textbf{Objetivos:} O1, O2 \quad \textbf{CE:} 3.a, 4.a \quad
\textbf{Duración estimada:} 10–15 minutos
\end{nota}

\paragraph{Descripción.} Juego de arrastrar y soltar en tres niveles, implementado con la API
\emph{Drag and Drop} de HTML5 (con alternativa táctil para móviles). Cada acierto suma 10
puntos; cada fallo resta 2 y muestra una explicación breve. Hay que superar un nivel para
desbloquear el siguiente.

\begin{reto}{Nivel 1 -- Clasificación de sentencias}
Arrastra cada sentencia a su sublenguaje: \textbf{DDL} (definición), \textbf{DML}
(manipulación), \textbf{DCL} (control de acceso) o \textbf{TCL} (control de transacciones).

\texttt{CREATE}, \texttt{SELECT}, \texttt{GRANT}, \texttt{COMMIT}, \texttt{ALTER},
\texttt{INSERT}, \texttt{REVOKE}, \texttt{ROLLBACK}, \texttt{DROP}, \texttt{UPDATE},
\texttt{TRUNCATE}, \texttt{DELETE}, \texttt{SAVEPOINT}.
\end{reto}

\begin{solucion}
\begin{tabularx}{\linewidth}{@{}>{\bfseries}lX@{}}
DDL & \texttt{CREATE}, \texttt{ALTER}, \texttt{DROP}, \texttt{TRUNCATE} \\
DML & \texttt{SELECT}, \texttt{INSERT}, \texttt{UPDATE}, \texttt{DELETE} \\
DCL & \texttt{GRANT}, \texttt{REVOKE} \\
TCL & \texttt{COMMIT}, \texttt{ROLLBACK}, \texttt{SAVEPOINT} \\
\end{tabularx}

\smallskip
\emph{Nota:} \texttt{TRUNCATE} se clasifica como DDL porque vacía la tabla
redefiniéndola, sin borrar fila a fila, y en muchos SGBD no admite \texttt{ROLLBACK}.
\end{solucion}

\begin{reto}{Nivel 2 -- Ordena las cláusulas}
Ordena las piezas para formar una consulta sintácticamente correcta que obtenga, para
cada socio con más de 400 kg entregados en 2026, el total de kilos, de mayor a menor:

\texttt{HAVING SUM(kilos) > 400} $\cdot$ \texttt{FROM entrega} $\cdot$
\texttt{ORDER BY total DESC} $\cdot$ \texttt{SELECT id\_socio, SUM(kilos) AS total} $\cdot$
\texttt{GROUP BY id\_socio} $\cdot$ \texttt{WHERE fecha >= '2026-01-01'}
\end{reto}

\begin{solucion}
\begin{lstlisting}
SELECT id_socio, SUM(kilos) AS total
FROM entrega
WHERE fecha >= '2026-01-01'
GROUP BY id_socio
HAVING SUM(kilos) > 400
ORDER BY total DESC;
\end{lstlisting}
Resultado: socio 1 (1050 kg), socio 2 (950 kg), socio 3 (700 kg). El socio 4
(350 kg) no supera el filtro del \texttt{HAVING}.

\smallskip
\emph{Pregunta adicional del nivel:} ordena las cláusulas según su orden \emph{lógico de
evaluación}. Solución: \texttt{FROM} $\rightarrow$ \texttt{WHERE} $\rightarrow$
\texttt{GROUP BY} $\rightarrow$ \texttt{HAVING} $\rightarrow$ \texttt{SELECT}
$\rightarrow$ \texttt{ORDER BY}.
\end{solucion}

\begin{reto}{Nivel 3 -- Relaciona operador y significado}
Une cada expresión con su significado.
\end{reto}

\begin{solucion}
\begin{tabularx}{\linewidth}{@{}lX@{}}
\toprule
\textbf{Expresión} & \textbf{Significado} \\
\midrule
\texttt{nombre LIKE 'Car\%'} & El nombre empieza por «Car» \\
\texttt{nombre LIKE '\_o\%'} & La segunda letra del nombre es una «o» \\
\texttt{municipio IN ('Lorca','Cieza')} & El municipio es Lorca o Cieza \\
\texttt{precio\_kg BETWEEN 1 AND 2} & Precio entre 1 y 2 €, ambos incluidos \\
\texttt{fecha\_alta IS NULL} & No se conoce la fecha de alta \\
\texttt{SELECT DISTINCT municipio} & Muestra cada municipio una sola vez \\
\texttt{SUM(kilos) AS total} & Pone el alias «total» a la columna calculada \\
\texttt{precio\_kg <> 0.60} & Precio distinto de 0,60 € \\
\bottomrule
\end{tabularx}
\end{solucion}

\begin{reto}{Llave del juego 1 -- dígito 1}
Al terminar los tres niveles se pregunta: \emph{¿Cuántas sentencias de la lista del
nivel~1 pertenecen al DCL?}
\end{reto}
\begin{solucion}
\textbf{2} (\texttt{GRANT} y \texttt{REVOKE}). \quad \textbf{Dígito 1 = 2}.
\end{solucion}

% ---------------------------------------------------------------------
\subsubsection{Juego 2: «¿Qué devuelve la consulta?»}
% ---------------------------------------------------------------------

\begin{nota}
\textbf{URL pública:} \url{\urlJuegoDos}\\
\textbf{Objetivos:} O2, O3, O4, O5, O6, O7 \quad
\textbf{CE:} 3.b, 3.c, 3.d, 3.e, 3.f, 4.b, 4.e, 4.h \quad
\textbf{Duración estimada:} 15–20 minutos
\end{nota}

\paragraph{Descripción.} Concurso de preguntas de respuesta múltiple con cuenta atrás
(30~s por pregunta) y puntuación según rapidez (entre 500 y 1000 puntos por acierto, como en
los concursos de televisión, más una bonificación de 100 puntos por racha a partir del tercer
acierto seguido). Cada pregunta muestra una sentencia SQL sobre la base de datos de la
\autoref{tab:datos}, que se puede consultar en todo momento con el botón «Ver tablas». Tras
responder, el juego indica si la respuesta es correcta y muestra el resultado real de la
consulta con una explicación. El docente puede proyectarlo en el aula para jugar por equipos
o el alumnado puede jugarlo de forma individual. Al final se muestra la clasificación
guardada en el navegador.

A continuación figuran las doce preguntas con sus opciones (la correcta se marca con
\textbf{\checkmark}).

\newcommand{\ok}{\textbf{\checkmark}}

\begin{reto}{Pregunta 1 -- Selección simple}
\begin{lstlisting}
SELECT nombre FROM socio WHERE municipio = 'Murcia';
\end{lstlisting}
a) Ana Martínez \quad
b) Ana Martínez y Pedro Ruiz \ok \quad
c) Pedro Ruiz \quad
d) Los cinco socios
\end{reto}

\begin{reto}{Pregunta 2 -- Función de agregado \emph{(llave: dígito 2)}}
\begin{lstlisting}
SELECT COUNT(*) FROM entrega;
\end{lstlisting}
a) 5 \quad b) 6 \quad c) 9 \ok \quad d) 3050
\end{reto}

\begin{reto}{Pregunta 3 -- Rango y ordenación}
\begin{lstlisting}
SELECT nombre FROM producto
WHERE precio_kg BETWEEN 1 AND 2
ORDER BY precio_kg DESC;
\end{lstlisting}
¿Qué producto aparece en la \emph{primera} fila?\\
a) Albaricoque \quad b) Alcachofa \ok \quad c) Pimiento \quad d) Melocotón
\end{reto}

\begin{reto}{Pregunta 4 -- Suma con filtro}
\begin{lstlisting}
SELECT SUM(kilos) FROM entrega WHERE id_socio = 1;
\end{lstlisting}
a) 500 \quad b) 700 \quad c) 1050 \ok \quad d) 3
\end{reto}

\begin{reto}{Pregunta 5 -- Agrupamiento y HAVING}
\begin{lstlisting}
SELECT id_socio, COUNT(*) FROM entrega
GROUP BY id_socio
HAVING COUNT(*) > 2;
\end{lstlisting}
¿Cuántas filas devuelve?\\
a) 0 \quad b) 1 \ok \quad c) 4 \quad d) 9
\end{reto}

\begin{reto}{Pregunta 6 -- Composición interna}
\begin{lstlisting}
SELECT s.nombre
FROM socio s
  INNER JOIN entrega e  ON s.id_socio = e.id_socio
  INNER JOIN producto p ON e.id_producto = p.id_producto
WHERE p.nombre = 'Alcachofa';
\end{lstlisting}
a) José López y Pedro Ruiz \ok \quad
b) Pedro Ruiz \quad
c) Ana Martínez y Carmen Sánchez \quad
d) Ninguna fila
\end{reto}

\begin{reto}{Pregunta 7 -- Composición externa}
\begin{lstlisting}
SELECT s.nombre
FROM socio s LEFT JOIN entrega e ON s.id_socio = e.id_socio
WHERE e.id_entrega IS NULL;
\end{lstlisting}
a) Ninguna fila \quad b) Lucía Gómez \ok \quad c) Los cinco socios \quad d) Error de sintaxis
\end{reto}

\begin{reto}{Pregunta 8 -- Subconsulta escalar}
\begin{lstlisting}
SELECT nombre FROM producto
WHERE precio_kg > (SELECT AVG(precio_kg) FROM producto);
\end{lstlisting}
a) Solo Albaricoque \quad
b) Alcachofa, Melocotón y Albaricoque \ok \quad
c) Limón y Naranja \quad
d) Pimiento, Alcachofa, Melocotón y Albaricoque
\end{reto}

\begin{reto}{Pregunta 9 -- Patrones con LIKE}
\begin{lstlisting}
SELECT COUNT(*) FROM socio WHERE nombre LIKE '%ez';
\end{lstlisting}
a) 2 \quad b) 3 \quad c) 4 \ok \quad d) 5
\end{reto}

\begin{reto}{Pregunta 10 -- Actualización}
\begin{lstlisting}
UPDATE producto SET precio_kg = precio_kg * 1.10
WHERE categoria = 'Cítrico';
\end{lstlisting}
¿Cuántas filas se modifican?\\
a) 0 \quad b) 1 \quad c) 2 \ok \quad d) 6
\end{reto}

\begin{reto}{Pregunta 11 -- Integridad referencial}
Si la clave ajena \texttt{entrega.id\_socio} se ha definido sin \texttt{ON DELETE CASCADE},
¿qué ocurre al ejecutar la siguiente sentencia?
\begin{lstlisting}
DELETE FROM socio WHERE id_socio = 1;
\end{lstlisting}
a) Se borra el socio y sus tres entregas \quad
b) Se borra el socio y sus entregas quedan con \texttt{id\_socio} a \texttt{NULL} \quad
c) El SGBD rechaza la sentencia por violar la integridad referencial \ok \quad
d) Se borran todos los socios
\end{reto}

\begin{reto}{Pregunta 12 -- Transacciones}
\begin{lstlisting}
START TRANSACTION;
DELETE FROM entrega WHERE kilos < 200;
ROLLBACK;
SELECT COUNT(*) FROM entrega;
\end{lstlisting}
a) 7 \quad b) 9 \ok \quad c) 0 \quad d) Error: no se puede borrar
\end{reto}

\begin{solucion}
\begin{tabularx}{\linewidth}{@{}clX@{}}
\toprule
\textbf{P.} & \textbf{Correcta} & \textbf{Justificación} \\
\midrule
1  & b & Solo Ana Martínez (id~1) y Pedro Ruiz (id~4) son de Murcia. \\
2  & c & La tabla \texttt{entrega} tiene 9 filas. \textbf{Dígito 2 = 9}. \\
3  & b & Están en el rango Alcachofa (1,80), Melocotón (1,50) y Pimiento (1,20); en orden descendente
         la primera es Alcachofa. Albaricoque (2,10) queda fuera. \\
4  & c & $500 + 200 + 350 = 1050$ kg. \\
5  & b & El socio 1 tiene 3 entregas; los socios 2, 3 y 4 tienen 2; el 5, ninguna. \\
6  & a & La alcachofa (id~4) aparece en las entregas 4 (socio 2) y 8 (socio 4). \\
7  & b & Lucía Gómez no tiene entregas: el \texttt{LEFT JOIN} rellena con \texttt{NULL}. \\
8  & b & Media $= 7{,}65/6 = 1{,}275$ €; la superan 1,80, 1,50 y 2,10. \\
9  & c & Martínez, López, Sánchez y Gómez acaban en «ez»; Ruiz, no. \\
10 & c & Hay dos cítricos: Limón y Naranja. \\
11 & c & Existen filas en \texttt{entrega} que referencian al socio 1 (restricción
         \texttt{RESTRICT}/\texttt{NO ACTION} por defecto). \\
12 & b & \texttt{ROLLBACK} deshace el borrado de las entregas~4 (150 kg) y~8 (100 kg),
         así que vuelve a haber 9 filas. \\
\bottomrule
\end{tabularx}
\end{solucion}

% ---------------------------------------------------------------------
\subsubsection{Juego 3: «Terminal de la Cooperativa»}
% ---------------------------------------------------------------------

\begin{nota}
\textbf{URL pública:} \url{\urlJuegoTres}\\
\textbf{Objetivos:} O3, O4, O5, O6, O7 \quad
\textbf{CE:} 3.c, 3.e, 3.f, 3.h, 4.b, 4.c, 4.d, 4.e, 4.f, 4.h \quad
\textbf{Duración estimada:} 30–40 minutos
\end{nota}

\paragraph{Descripción.} Es el reto más exigente y el más parecido al trabajo real.
La página simula la consola del servidor recuperado: incluye un editor de texto con resaltado
de sintaxis y un motor SQLite compilado a WebAssembly (biblioteca \emph{sql.js}), que carga la
base de datos \texttt{cooperativa} en el navegador. Así el alumnado escribe y ejecuta
\emph{sentencias SQL reales} sin instalar nada y sin depender de ningún servidor.

El juego plantea seis «incidencias» de la cooperativa. En cada una, el alumnado escribe la
sentencia y pulsa \emph{Comprobar}. El juego ejecuta la sentencia, compara el resultado
(o el estado de la tabla tras una modificación) con el esperado e informa del acierto o del
error. Cada incidencia vale 500 puntos. Si se equivoca, puede pedir hasta dos pistas, que restan 100 puntos cada una, y cada comprobación fallida resta 25, con un mínimo de 150 puntos por incidencia resuelta. En las primeras incidencias
se ofrecen plantillas con huecos para completar (modo guiado). En las últimas, el editor está
vacío (modo libre). Un botón \emph{Reiniciar base de datos} restaura los datos originales.
Al terminar, el juego muestra el resumen de puntos y la insignia «Maestro del DML».

\begin{reto}{Incidencia 1 -- Liquidación por socio (modo guiado)}
La gerencia necesita el importe que corresponde a cada socio (kilos $\times$ precio),
ordenado de mayor a menor. Completa los huecos:
\begin{lstlisting}
SELECT s.nombre, SUM(e.kilos * p.precio_kg) AS importe
FROM socio s
  ______ JOIN entrega e  ON s.id_socio = e.______
  INNER JOIN producto p ON e.id_producto = p.id_producto
______ s.nombre
ORDER BY importe ______;
\end{lstlisting}
\end{reto}
\begin{solucion}
Huecos: \texttt{INNER}, \texttt{id\_socio}, \texttt{GROUP BY}, \texttt{DESC}.

\begin{tabular}{@{}lr@{}}
\toprule
\textbf{nombre} & \textbf{importe} \\
\midrule
Ana Martínez   & 1065,00 \\
Carmen Sánchez &  780,00 \\
José López     &  630,00 \\
Pedro Ruiz     &  480,00 \\
\bottomrule
\end{tabular}
\end{solucion}

\begin{reto}{Incidencia 2 -- Categorías estrella (modo guiado) \emph{(llave: dígito 4)}}
Obtén las categorías de producto cuyo total de kilos entregados supere los 1000 kg.
\begin{lstlisting}
SELECT p.categoria, SUM(e.kilos) AS total_kg
FROM producto p INNER JOIN entrega e ON p.id_producto = e.id_producto
GROUP BY ______
______ SUM(e.kilos) > 1000;
\end{lstlisting}
\end{reto}
\begin{solucion}
Huecos: \texttt{p.categoria}, \texttt{HAVING}.
Resultado: una única fila, \emph{Cítrico -- 1600 kg} (Hortaliza suma 700 kg y Fruta de
hueso, 750 kg). \textbf{Dígito 4 = 1} (número de categorías obtenidas).
\end{solucion}

\begin{reto}{Incidencia 3 -- Producto sin movimiento (modo libre)}
Escribe una consulta que muestre los productos de los que no hay ninguna entrega.
Después, escribe una segunda versión con otra técnica y explica cuál prefieres.
\end{reto}
\begin{solucion}
\begin{lstlisting}
-- Versión con subconsulta
SELECT nombre FROM producto
WHERE id_producto NOT IN (SELECT id_producto FROM entrega);

-- Versión con composición externa
SELECT p.nombre
FROM producto p LEFT JOIN entrega e ON p.id_producto = e.id_producto
WHERE e.id_entrega IS NULL;
\end{lstlisting}
Resultado en ambos casos: \emph{Albaricoque}. Las dos versiones son válidas. La versión con
subconsulta se lee mejor, pero \texttt{NOT IN} devuelve un conjunto vacío si la subconsulta
contiene algún \texttt{NULL}. La composición externa (o \texttt{NOT EXISTS}) evita ese
problema y los optimizadores suelen ejecutarla de forma eficiente (CE~3.h).
\end{solucion}

\begin{reto}{Incidencia 4 -- Nuevo socio (modo libre)}
Se incorpora a la cooperativa Antonio Pérez, de Yecla, con fecha de alta 1 de octubre de
2026 e identificador 6. Regístralo y, a continuación, registra que hoy (2026-10-29) ha
entregado 600 kg de limones (entrega número 10).
\end{reto}
\begin{solucion}
\begin{lstlisting}
INSERT INTO socio (id_socio, nombre, municipio, fecha_alta)
VALUES (6, 'Antonio Pérez', 'Yecla', '2026-10-01');

INSERT INTO entrega (id_entrega, id_socio, id_producto, fecha, kilos)
VALUES (10, 6, 1, '2026-10-29', 600);
\end{lstlisting}
El orden importa: si se inserta primero la entrega, el SGBD la rechaza porque el socio
6 aún no existe (integridad referencial, CE~4.h).
\end{solucion}

\begin{reto}{Incidencia 5 -- Limpieza de entregas antiguas (modo libre) \emph{(llave: dígito 3)}}
\emph{Partiendo de la base de datos original (pulsa Reiniciar).} Dentro de una
transacción, borra las entregas anteriores al 6 de marzo de 2026, comprueba cuántas filas
se han borrado y confirma los cambios.
\end{reto}
\begin{solucion}
\begin{lstlisting}
START TRANSACTION;          -- BEGIN TRANSACTION en SQLite
DELETE FROM entrega WHERE fecha < '2026-03-06';
SELECT COUNT(*) FROM entrega;   -- quedan 6 filas
COMMIT;
\end{lstlisting}
Se borran \textbf{3 filas} (entregas 1, 2 y 3). \textbf{Dígito 3 = 3}.
\end{solucion}

\begin{reto}{Incidencia 6 -- Subida de precios con marcha atrás (modo libre)}
Dentro de una transacción: sube un 10\,\% el precio de las hortalizas, crea un punto de
guardado, sube otro 5\,\% el de la fruta de hueso, deshaz solo esta segunda subida y
confirma la primera. Por último, guarda en una tabla nueva \texttt{resumen\_categoria}
el precio medio de cada categoría.
\end{reto}
\begin{solucion}
\begin{lstlisting}
START TRANSACTION;
UPDATE producto SET precio_kg = precio_kg * 1.10
WHERE categoria = 'Hortaliza';
SAVEPOINT tras_hortalizas;
UPDATE producto SET precio_kg = precio_kg * 1.05
WHERE categoria = 'Fruta de hueso';
ROLLBACK TO SAVEPOINT tras_hortalizas;
COMMIT;

CREATE TABLE resumen_categoria (
  categoria    VARCHAR(30) PRIMARY KEY,
  precio_medio DECIMAL(5,2)
);
INSERT INTO resumen_categoria (categoria, precio_medio)
SELECT categoria, ROUND(AVG(precio_kg), 2)
FROM producto
GROUP BY categoria;
\end{lstlisting}
Resultado esperado: Pimiento pasa a 1,32 € y Alcachofa a 1,98 €; la fruta de hueso
conserva sus precios. \texttt{resumen\_categoria}: Cítrico 0,53; Fruta de hueso 1,80;
Hortaliza 1,65 (CE~4.c, 4.d, 4.f).
\end{solucion}

\subsection{Código de la caja fuerte}

\begin{center}
\begin{tabular}{@{}lll c@{}}
\toprule
\textbf{Dígito} & \textbf{Juego} & \textbf{Reto} & \textbf{Valor} \\
\midrule
1 & Juego 1 & Sentencias DCL de la lista & 2 \\
2 & Juego 2 & Pregunta 2: \texttt{COUNT(*)} de \texttt{entrega} & 9 \\
3 & Juego 3 & Incidencia 5: filas borradas & 3 \\
4 & Juego 3 & Incidencia 2: categorías con más de 1000 kg & 1 \\
\midrule
\multicolumn{3}{r}{\textbf{Código de la caja fuerte}} & \textbf{\Large 2931} \\
\bottomrule
\end{tabular}
\end{center}

\subsection{Secuenciación y metodología}

La actividad está pensada para tres sesiones de 55 minutos, en el momento de la programación
en que se empiezan a trabajar los RA3 y RA4:

\begin{itemize}
  \item \textbf{Sesión 1}: presentación de la misión y de la base de datos, introducción a SQL,
        \textbf{juego~1}, consultas simples y de resumen.
  \item \textbf{Sesión 2}: composiciones, subconsultas y \textbf{juego~2} jugado en el aula
        por equipos con clasificación proyectada.
  \item \textbf{Sesión 3}: manipulación de datos y transacciones, \textbf{juego~3} por
        parejas y apertura de la caja fuerte. Puesta en común y autoevaluación.
\end{itemize}

Se usa una metodología activa que combina la explicación interactiva con la
\emph{presentación}, el aprendizaje basado en retos y el trabajo cooperativo. Como los
juegos son páginas web abiertas, el alumnado puede repetirlos fuera del aula tantas veces
como quiera, lo que facilita la atención a la diversidad: quien necesita más práctica puede
volver a jugarlos, y quien avanza rápido puede abordar en el juego~3 las versiones
alternativas de cada consulta. Las páginas siguen pautas básicas de accesibilidad
(contraste suficiente, navegación por teclado y textos alternativos).

\paragraph{Evaluación.} Los juegos 1 y 2 tienen carácter formativo: dan
retroalimentación inmediata al alumnado y permiten al docente detectar errores
frecuentes. Las soluciones del juego~3 (que el alumnado puede copiar o exportar
desde la propia terminal) sirven como evidencia para evaluar los criterios 3.c–3.f, 3.h y
4.b–4.f con una rúbrica.

% =====================================================================
\section{Bibliografía}
% =====================================================================

\nocite{*}
\printbibliography[heading=none]

% =====================================================================
\appendix
\section{Guion SQL de la base de datos \texttt{cooperativa}}
\label{ap:script}
% =====================================================================

El siguiente guion crea y carga la base de datos de trabajo. Es compatible con SQLite,
MySQL/MariaDB y PostgreSQL. Es el mismo que el juego~3 carga en el navegador y está
disponible en el repositorio de los juegos.

\lstinputlisting{../web/data/cooperativa.sql}

\end{document}
